\chapter{Irritable Bowel Syndrome}
\index{Irritable Bowel Syndrome}
\label{sec:Irritable Bowel Syndrome}

\section{Overview}
Irritable bowel syndrome (IBS) is a chronic functional digestive tract disorder characterized by recurrent abdominal pain associated with altered bowel habits, in the absence of structural or biochemical abnormalities. It affects approximately 10--15\% of the global population and is more common in women. This genetic disorder results from specific gene mutations or chromosomal abnormalities. It may be present at birth or manifest later in life depending on the underlying mechanism. Genetic disorders result from abnormalities in an individual's DNA, ranging from single-gene mutations to chromosomal abnormalities. These conditions may be inherited from parents or arise from new mutations. Pattern of inheritance depends on whether the mutation is dominant, recessive, or X-linked. Genetic counseling helps families understand risks and make informed decisions.
\section{Causes}
This condition happens because of visceral hypersensitivity, altered digestive tract motility, gut-brain axis dysfunction, intestinal dysbiosis, low-grade mucosal inflammation, and post-infectious triggers. IBS is subtyped by predominant stool pattern: IBS with constipation (IBS-C), IBS with diarrhea (IBS-D), IBS with mixed bowel habits (IBS-M), and unsubtyped IBS (IBS-U). Genetic disorders result from mutations in specific genes that encode proteins essential for normal cellular function. The pattern of inheritance depends on whether the mutation is dominant, recessive, or X-linked.   
\section{Risk Factors}

\begin{itemize}[noitemsep]
  \item Genetic predisposition and family history
  \item Advancing age
  \item Lifestyle factors (diet, exercise, smoking, alcohol)
  \item Environmental exposures
  \item Underlying medical conditions
  \item Medication use
\end{itemize}
\section{Signs and Symptoms}

\subsection{Common symptoms}
\begin{itemize}[noitemsep]
  \item \item You may experience recurrent abdominal pain, on average at least 1 day per week in the last 3 months, associated with defecation, change in stool frequency, or change in stool form.

  \item Pain or discomfort in the affected area
  \end{itemize}

\subsection{Emergency warning signs}
\begin{itemize}[noitemsep]
  \item Severe or worsening symptoms of irritable bowel syndrome require immediate medical evaluation
\end{itemize}
\section{Progression and Stages}
The condition may progress through different stages of severity over time. Early stages often involve mild or subtle symptoms that may be overlooked. Moderate stages involve more noticeable symptoms affecting daily function. Advanced stages may involve significant impairment and complications.
\section{Complications}
\begin{itemize}[noitemsep]
  \item Disease progression leading to worsening symptoms
  \item Secondary infections or injuries
  \item Side effects of treatment
  \item Reduced quality of life
  \item Emotional and psychological impact
\end{itemize}
\section{Diagnosis}
\begin{itemize}[noitemsep]
  \item Doctors diagnose this based on the Rome IV criteria
  \item Alarm features (rectal bleeding, nocturnal symptoms, unintentional weight loss, family history of colorectal cancer, iron deficiency anemia, onset after age 50) warrant further investigation to exclude organic disease.
  \item Comprehensive medical history and physical examination
  \item Laboratory tests to assess organ function
  \item Imaging studies to visualize affected structures
  \item Specialized testing depending on the affected system
  \item Consultation with relevant specialists
\end{itemize}
\section{Prevention}
\begin{itemize}[noitemsep]
  \item Maintain a healthy lifestyle with balanced diet and exercise
  \item Avoid known risk factors
  \item Attend regular medical check-ups for early detection
  \item Follow recommended screening guidelines
\end{itemize}
\section{Treatment Options}
Treatment typically involves multimodal: dietary modification (low FODMAP diet, soluble fiber supplementation), antispasmodics (hyoscine, dicyclomine, peppermint oil), laxatives for IBS-C (linaclotide, lubiprostone, plecanatide), antidiarrheals for IBS-D (loperamide, eluxadoline, rifaximin), neuromodulators (tricyclic antidepressants, SSRIs), and gut-directed psychological therapies (hypnotherapy, thinking and memory-behavioral therapy). Symptomatic management and supportive care Enzyme replacement therapy for certain conditions Gene therapy (emerging for select conditions) Dietary modifications (e.g., phenylketonuria) Regular monitoring for associated complications Physical and occupational therapy Genetic counseling for family planning Participation in clinical trials for new therapies

\begin{itemize}[noitemsep]
  \item Medications to manage symptoms and address underlying mechanisms
  \item Lifestyle modifications and self-care measures
  \item Physical therapy and rehabilitation
  \item Surgical intervention when indicated
  \item Regular monitoring and follow-up care
\end{itemize}
\section{Living with It}
\begin{itemize}[noitemsep]
  \item Follow your treatment plan as prescribed
  \item Maintain regular follow-up with your healthcare provider
  \item Make necessary lifestyle adjustments
  \item Seek support from family, friends, or support groups
  \item Stay informed about your condition
\end{itemize}
\section{Outlook}
\begin{itemize}[noitemsep]
  \item The outlook varies from person to person
  \item Symptoms are chronic but can be managed with multimodal therapy.
\end{itemize}
